\section{Futuristic CAD and physical-transfer revision}
All 122 actual futuristic FreeCAD surfaces are imported into an explicit provisional six-joint assembly. The CAD hip, knee and ankle landmarks define its proportions and bent neutral pose. Estimated mass is 0.500 kg; actuator dynamics, mass properties, mechanical clearance and self-collision are not hardware-validated.

The earlier long run was stopped after validation and interactive use exposed poor backward heading, rough motion and transition failures. Good planar speed error did not establish command reliability. The replacement revision uses 41 observations: the earlier 29 plus applied-target and target-velocity state. A shared CPU/GPU controller ramps commands, filters complete motor targets, and limits target speed. A 1.5 Hz reference cadence with 8 mm lift replaces the faster reference. Extra yaw-error and complete-target velocity/acceleration costs address the reported failures. Left/right mirror consistency uses the maintained RSL-RL extension. Prior actor/critic mappings are retained by zero-padding the new input channels, with a tested fresh-optimizer weight transfer.

The new sensor-to-position runtime exactly matches simulator targets in a 100-step handoff check. The exported ONNX runtime loads without MuJoCo, mjlab or RSL-RL; it performs no hardware I/O. The 21 software checks include randomized GPU execution, reset isolation, target-rate bounds, observation migration, and mirror/gait consistency. Passing these checks is not walking acceptance.

The randomized 20,000-update training stage is in progress on the local RTX 4060. It varies physically consistent mass/inertia, friction, gains, effort, actuation delay, and gyro/encoder bias and noise. These ranges model uncertainty and are not calibrated servo measurements. Validation prioritizes the weakest command group and worst yaw error; held-out physical-parameter draws and CPU command transitions remain separate gates. The revised policy is not yet qualified for all commands. Actual servo identities, mass/COM, sensor feedback and timing are still needed; physical hardware performance remains unmeasured. The final artifact-linked tables will replace this progress section after the scheduled validation/test pipeline completes.
